\documentclass[11pt]{article}
\usepackage{amsmath,amsthm,amssymb}
\usepackage[margin=1in]{geometry}
\usepackage{enumitem}
\usepackage{booktabs}

\newtheorem{theorem}{Theorem}
\newtheorem{lemma}[theorem]{Lemma}
\newtheorem{proposition}[theorem]{Proposition}
\newtheorem{corollary}[theorem]{Corollary}
\newtheorem{conjecture}[theorem]{Conjecture}
\theoremstyle{definition}
\newtheorem{definition}[theorem]{Definition}
\newtheorem{remark}[theorem]{Remark}
\newtheorem{example}[theorem]{Example}

\title{Closing the hook rank formula:\\
the property $(Q_m)$ via strong image and sign induction}
\author{Clio}
\date{10 May 2026}

\begin{document}
\maketitle

\begin{abstract}
The May~9 writeup reduced the hook rank formula
$\mathrm{rank}\,\Pi^{S_n}|_{V_{(n-1,1)}}=\lceil(n-2)/2\rceil$ to the
single property
\[ (Q_m): \quad v_{m-1}^*(K_m) \neq 0 \quad \text{for all even } m \ge 4, \]
where $K_m := \ker\Pi^{S_m}|_{V_{(m-1,1)}}$.  Here we prove $(Q_m)$
unconditionally, completing the program.

The proof rests on two pieces.  First, a \textbf{Strong Image
Coefficient Formula}: explicit triangular inversion of the operator
$R_{n-1}'' = (T_{n-2}{+}1)\cdots(T_1{+}1)$ on $V_{(n-1,1)}$ shows that
for any $w$ with $R_n'(w) = y \in \mathrm{span}(v_2,\ldots,v_{n-2})$,
\[
   w_{v_l} \;=\; \frac{y_{v_{l-1}} - \gamma_{l-2}\, y_{v_{l-2}}}{(1+q)^{n-2}},
   \qquad 3 \le l \le n-1,
\]
where $\gamma_i := \beta_i p_{i+1}$ are explicit Hoefsmit
constants. Second, an alternating \textbf{sign induction} on the
``slopes'' $\mu_l$ (for $l$ odd) and $\nu_l$ (for $l$ even) on the
new-generator quotient of $K_l$.  At any positive real $q$, both
slopes are strictly negative, while $\gamma_{m-3}$ is positive --
forcing $\mu_{m-1} - \gamma_{m-3} < 0$, equivalently $(Q_m)$.

Combined with the May-9 reduction theorem, this gives the hook rank
formula in full unconditionally.
\end{abstract}

\section{Setup}\label{sec:setup}

We retain all conventions from the May-8 morning, May-8 evening, and
May-9 papers (\texttt{2026-05-08-hook-rank-formula.tex},
\texttt{2026-05-08-evening-2-kernel-dichotomy.tex},
\texttt{2026-05-09-nesting-and-reduction.tex}). In particular:
\begin{itemize}[topsep=0.3em,itemsep=0.2em,leftmargin=2em]
\item $V_{(n-1,1)}$ has the seminormal basis $v_2,\ldots,v_n$.
\item $u_i = v_i + \beta_i v_{i+1}$ is the $+q$-eigenvector of $T_i$
  in our normalization $\alpha_i = 1, q_i = 1/(1+q)$, with
  \[
     \beta_i \;=\; \frac{q[i-1]}{[i]}, \qquad
     \beta'_i \;=\; -\frac{[i+1]}{[i]}, \qquad
     p_i \;=\; \frac{[i+1]}{[2][i]}, \qquad
     q_i \;=\; \frac{1}{[2]},
  \]
  where $[k] := (q^k-1)/(q-1) = 1 + q + \cdots + q^{k-1}$.
\item $K_m := \ker\Pi^{S_m}|_{V_{(m-1,1)}}$, $D_m := R_m'(V_{(m-1,1)})
  = \mathrm{span}(v_2,\ldots,v_{m-2}, u_{m-1})$.
\item $(P_m)$: $v_m^*(K_m) = 0$ iff $m$ odd.
\item $(Q_m)$: $v_{m-1}^*(K_m) \neq 0$.
\item $(R_m)$: for $m$ even, $v_{m-1}^*$ and $v_m^*$ are proportional
  on $K_m$ with nonzero ratio.
\item Image Lemma (May-8 evening): for $w \in V_{(m-1,1)}$ with
  $R_m'(w) = y \in \mathrm{span}(v_2,\ldots,v_{m-2})$, $w_{v_m} = C_m
  y_{v_{m-2}}$ with $C_m = \beta_{m-2}\beta'_{m-1}/(1+q)^{m-1}$.
\item Nesting (May-9): for all $n \ge 4$, $K_{n-2}^{\mathrm{embed}}
  \subset K_n$ via the natural embedding.
\end{itemize}

We define
\[
   \gamma_i \;:=\; \beta_i p_{i+1} \;=\; \frac{q[i-1][i+2]}{[2][i][i+1]}
   \qquad (i \ge 2),
\]
which is positive at every positive real $q$.

\section{The Strong Image Coefficient Formula}\label{sec:strong-image}

\begin{theorem}[Strong Image Coefficient Formula]\label{thm:strong-image}
Let $n \ge 4$ and let $w \in V_{(n-1,1)}$ satisfy
$R_n'(w) = y \in \mathrm{span}(v_2,\ldots,v_{n-2})$.  Modulo
$\ker R_n' = \mathbb{Q}(q) v_2$, the coordinate $w_{v_l}$ is uniquely
determined for $3 \le l \le n$, and
\begin{equation}\label{eq:strong-image}
   w_{v_l} \;=\; \frac{y_{v_{l-1}} - \gamma_{l-2}\, y_{v_{l-2}}}{(1+q)^{n-2}},
   \qquad 3 \le l \le n-1,
\end{equation}
with the convention $\gamma_1 = 0, y_{v_1} = 0$.  Furthermore,
$w_{v_n} = C_n y_{v_{n-2}}$ (the May-8 evening Image Lemma).
\end{theorem}

\begin{proof}
Write $R_n' = (T_{n-1}{+}1) R_{n-1}''$ with
$R_{n-1}'' = (T_{n-2}{+}1)\cdots(T_1{+}1)$.  The image of $R_{n-1}''$
is $I_{n-2} = \mathrm{span}(v_2,\ldots,v_{n-3}, u_{n-2}, v_n)$
(Lemma~4 of the May-8 morning paper).  The map
$R_{n-1}''\colon V_{(n-1,1)} \to I_{n-2}$ has 1-dim kernel
$\mathbb{Q}(q) v_2$ since $(T_1{+}1) v_2 = 0$.

Set $z := R_{n-1}''(w) \in I_{n-2}$.  We claim that the matrix of
$R_{n-1}''$ from the source basis $(v_3, v_4, \ldots, v_n)$ (in the
quotient $V_{(n-1,1)}/\mathbb{Q}(q) v_2$) to the image basis
$(v_2, v_3, \ldots, v_{n-3}, u_{n-2}, v_n)$ is lower-triangular with
constant nonzero diagonal $(1+q)^{n-3}$.

\medskip
\textbf{Computation of the matrix.}  Let
$L_j := (T_j{+}1)(T_{j-1}{+}1)\cdots(T_1{+}1)$.  We track $L_j v_l$
for $l \ge 3$ and $j = 1, 2, \ldots, n-2$.  By induction on $j$,
\begin{equation}\label{eq:Ljvl}
   L_j v_l \;=\;
   \begin{cases}
   (1+q)^j q_{l-1} \bigl(v_{l-1} + B_{l,l} v_l + B_{l,l+1} v_{l+1} +
   \cdots + B_{l,j-1} v_{j-1} + B_{l,j} u_j\bigr) & \text{if } 3 \le l \le j+1,\\
   (1+q)^j v_l & \text{if } l \ge j+2,
   \end{cases}
\end{equation}
where $B_{l,k} := \prod_{i=l-1}^{k-1}\beta_i p_{i+1} = \prod_{i=l-1}^{k-1}\gamma_i$
(empty product equals $1$).

The base $j = 1$ gives $L_1 v_l = (1+q) v_l$ for $l \ge 3$.  For the
inductive step $j-1 \to j$: the operator $T_j{+}1$ acts as $(1+q)$ on
$v_k$ for $k \ne j, j+1$, and on the block $\mathrm{span}(v_j, v_{j+1})$ as
\[
   (T_j{+}1) v_j = (1+q) p_j u_j,
   \qquad
   (T_j{+}1) v_{j+1} = u_j
   = v_j + \beta_j v_{j+1},
\]
using $q_j = 1/(1+q)$.  Thus $(T_j{+}1) u_{j-1} = (T_j{+}1)(v_{j-1} +
\beta_{j-1} v_j) = (1+q) v_{j-1} + (1+q) \beta_{j-1} p_j u_j =
(1+q)(v_{j-1} + \gamma_{j-1} u_j)$.  Substituting these into the
inductive form for $L_{j-1}$ and collecting yields \eqref{eq:Ljvl}.

\medskip
\textbf{Reading off the matrix.} Apply at $j = n-2$.  For source
$v_l$ with $3 \le l \le n-1$, the image $L_{n-2} v_l$ is supported on
the image-basis vectors $v_{l-1}, v_l, \ldots, v_{n-3}, u_{n-2}$,
with leading coefficient $(1+q)^{n-2} q_{l-1} = (1+q)^{n-3}$ (using
$q_{l-1} = 1/(1+q)$) at $v_{l-1}$.  For $l = n-1$: $L_{n-2}
v_{n-1} = (1+q)^{n-2} q_{n-2} u_{n-2} = (1+q)^{n-3} u_{n-2}$.  For
$l = n$: $L_{n-2} v_n = (1+q)^{n-2} v_n$.

In short, the matrix $M$ from source basis $(v_3, v_4, \ldots, v_n)$
to image basis $(v_2, v_3, \ldots, v_{n-3}, u_{n-2}, v_n)$ has:
\begin{itemize}[topsep=0.3em,itemsep=0.2em,leftmargin=2em]
\item Column for $v_l$, $3 \le l \le n-1$: leading $(1+q)^{n-3}$ at
  row $v_{l-1}$ (or row $u_{n-2}$ if $l = n-1$); off-diagonal entries
  $(1+q)^{n-3} \prod_{i=l-1}^{k-1}\gamma_i$ at row $v_k$ (or $u_{n-2}$
  for $k = n-2$), for $l \le k \le n-2$.
\item Column for $v_n$: $(1+q)^{n-2}$ at row $v_n$, zero elsewhere.
\end{itemize}
Reading column-by-column, this is lower-triangular with all diagonal
entries equal to $(1+q)^{n-3}$ in the rows above $v_n$, and
$(1+q)^{n-2}$ on the $v_n$-diagonal.

\medskip
\textbf{Back-substitution.}  Solve $z = M w$ row by row.  Index the
source basis as $w_3, w_4, \ldots, w_n$ and the image basis as
$z^{(v_2)}, z^{(v_3)}, \ldots, z^{(v_{n-3})}, z^{(u_{n-2})}, z^{(v_n)}$.

\emph{Row $v_2$:} $z^{(v_2)} = (1+q)^{n-3} w_3$, so
$w_3 = z^{(v_2)}/(1+q)^{n-3}$.

\emph{Row $v_k$, $3 \le k \le n-3$:}
$z^{(v_k)} = (1+q)^{n-3}\!\left[\sum_{l=3}^{k}\Pi_{l,k}\, w_l \,+\, w_{k+1}\right]$,
where $\Pi_{l,k} := \prod_{i=l-1}^{k-1}\gamma_i$.

We claim by induction on $k$ that
\begin{equation}\label{eq:wk-formula}
   w_{k+1} \;=\; \frac{z^{(v_k)} - \gamma_{k-1}\, z^{(v_{k-1})}}{(1+q)^{n-3}}
\end{equation}
for $3 \le k \le n-3$ (with $z^{(v_2)}$ taking the place of
$z^{(v_{k-1})}$ at $k = 3$).  The case $k = 3$:
$z^{(v_3)} = (1+q)^{n-3}[\gamma_2 w_3 + w_4]$, so
$w_4 = z^{(v_3)}/(1+q)^{n-3} - \gamma_2 w_3
     = (z^{(v_3)} - \gamma_2 z^{(v_2)})/(1+q)^{n-3}$, matching
\eqref{eq:wk-formula}.

For the inductive step $k-1 \to k$: by the inductive hypothesis,
$w_{l+1} = (z^{(v_l)} - \gamma_{l-1} z^{(v_{l-1})})/(1+q)^{n-3}$ for
$l < k$.  Substituting into the row-$v_k$ equation and using the
telescoping identity $\Pi_{l,k} = \gamma_{l-1}\Pi_{l-1,k}$ -- which
follows directly from $\Pi_{l,k} = \prod_{i=l-1}^{k-1}\gamma_i$ --
all terms with $z^{(v_l)}$ for $l < k-1$ cancel, leaving exactly
\eqref{eq:wk-formula}.  Specifically, after substituting, the
contribution of $z^{(v_l)}$ to $w_{k+1}$ is
\[
\Bigl[-\Pi_{l+1,k} + \Pi_{l+2,k} \gamma_l\Bigr]
\cdot \frac{z^{(v_l)}}{(1+q)^{n-3}}
\;=\;
\Bigl[-\Pi_{l+1,k} + \Pi_{l+1,k}\Bigr]\cdot \frac{z^{(v_l)}}{(1+q)^{n-3}} = 0
\]
for $2 \le l \le k-2$ (using $\Pi_{l+2,k}\gamma_l = \Pi_{l+1,k}$),
while the $z^{(v_{k-1})}$ contribution survives as
$-\gamma_{k-1}\, z^{(v_{k-1})}/(1+q)^{n-3}$.

\emph{Row $u_{n-2}$:} $z^{(u_{n-2})} =
(1+q)^{n-3}\bigl[\sum_{l=3}^{n-2}\Pi_{l,n-2} w_l + w_{n-1}\bigr]$.
By an identical telescoping argument, all $z^{(v_l)}$ contributions
for $l \le n-4$ cancel, leaving
\begin{equation}\label{eq:wnm1-formula}
   w_{n-1} \;=\; \frac{z^{(u_{n-2})} - \gamma_{n-3}\, z^{(v_{n-3})}}{(1+q)^{n-3}}.
\end{equation}

\medskip
\textbf{Translating to $y$.}  The element $y \in
\mathrm{span}(v_2,\ldots,v_{n-2})$ relates to $z$ via the constraint
$(T_{n-1}{+}1) z = y$.  Since $y$ has no $v_{n-1}$ or $v_n$
component, and $T_{n-1}{+}1$ has eigenvalues $1+q$ on
$\mathrm{span}(v_2,\ldots,v_{n-2})$ and $0$ on $\mathrm{span}(u'_{n-1})$:
\[
   z \;=\; \tfrac{1}{1+q} y \;+\; \mu\, u'_{n-1}
\]
for some $\mu \in \mathbb{Q}(q)$ chosen so that $z \in I_{n-2}$.
Working out the constraint (May-8 evening, eq.~(2.2)) yields
\[
   z^{(v_l)} = y_{v_l}/(1+q),\quad 2 \le l \le n-3, \qquad
   z^{(u_{n-2})} = y_{v_{n-2}}/((1+q)\alpha_{n-2}) = y_{v_{n-2}}/(1+q),
\]
using $\alpha_{n-2} = 1$.

Substituting into \eqref{eq:wk-formula} and \eqref{eq:wnm1-formula}:
for $3 \le k \le n-3$,
\[
   w_{k+1} = \frac{y_{v_k}/(1+q) - \gamma_{k-1} y_{v_{k-1}}/(1+q)}{(1+q)^{n-3}}
           = \frac{y_{v_k} - \gamma_{k-1}\, y_{v_{k-1}}}{(1+q)^{n-2}};
\]
and for $w_{n-1}$:
\[
   w_{n-1} = \frac{y_{v_{n-2}}/(1+q) - \gamma_{n-3} y_{v_{n-3}}/(1+q)}{(1+q)^{n-3}}
           = \frac{y_{v_{n-2}} - \gamma_{n-3}\, y_{v_{n-3}}}{(1+q)^{n-2}}.
\]
Setting $l = k+1$ in the first and matching the second: \eqref{eq:strong-image} is established for all $3 \le l \le n-1$.

The case $l = n$ (Image Lemma) was proved May-8 evening.
\end{proof}

\begin{remark}\label{rem:scope}
The Strong Image formula \eqref{eq:strong-image} extends the Image
Lemma from $w_{v_n}$ (which depends only on $y_{v_{n-2}}$) to every
``intermediate'' coordinate $w_{v_l}$ for $l = 3,\ldots,n-1$, each of
which depends on \emph{exactly two} consecutive coordinates of $y$.

This local structure is the key technical ingredient in what
follows: it converts the question of $(Q_m)$ on $K_m$ into a
linear-functional non-vanishing question on $K_{m-1}$.
\end{remark}

\section{Slopes on the new generators}\label{sec:slopes}

We now use the Strong Image formula to extract two ``slopes'' on the
new-generator quotient of $K_l$, indexed by parity.

\begin{definition}[Slopes $\mu_l$ and $\nu_l$]\label{def:slopes}
For $l \ge 4$, let $\xi^{(l)} \in K_l$ be any vector with
$\xi^{(l)} \notin K_{l-2}^{\mathrm{embed}}$ (such a vector exists
since $K_{l-2}^{\mathrm{embed}}$ has codimension $1$ in $K_l$).
\begin{itemize}[topsep=0.3em,itemsep=0.2em,leftmargin=2em]
\item For $l$ \textbf{odd} ($l \ge 5$): $\xi^{(l)}_{v_l} = 0$ by
  $(P_l)$ and $\xi^{(l)}_{v_{l-1}} \neq 0$ by $(Q_l)$ (the latter is
  automatic from $(P_l), (P_{l-1})$ as shown in May-9 \S2). Define
  \[
     \mu_l \;:=\; \frac{\xi^{(l)}_{v_{l-1}}}{\xi^{(l)}_{v_{l-2}}}
     \quad (\text{when } \xi^{(l)}_{v_{l-2}} \neq 0).
  \]
\item For $l$ \textbf{even} ($l \ge 4$): $\xi^{(l)}_{v_l} \neq 0$ by
  $(P_l)$.  Define
  \[
     \nu_l \;:=\; \frac{\xi^{(l)}_{v_{l-1}}}{\xi^{(l)}_{v_l}}.
  \]
\end{itemize}
Both are independent of the choice of $\xi^{(l)}$ within its coset
modulo $K_{l-2}^{\mathrm{embed}}$ (since $K_{l-2}^{\mathrm{embed}}
\subset \mathrm{span}(v_2,\ldots,v_{l-2})$ has zero $v_{l-1}, v_l$
components).
\end{definition}

\begin{lemma}[Base values]\label{lem:base-slopes}
$\nu_4 = -1/\gamma_2$ and $\mu_5 = -\gamma_2$.
\end{lemma}

\begin{proof}
$K_4 = \mathrm{span}(v_2, v_3 - c v_4)$ with $c = \gamma_2$ (verified
in May-9 nesting paper, via $c = q_2 \beta_2 p_3/q_3$).  The new
generator $\eta^{(4)} = v_3 - \gamma_2 v_4$ gives $\nu_4 = 1/(-\gamma_2) =
-1/\gamma_2$.

For $\mu_5$: $K_5 = \mathrm{span}(v_2, v_3 - c v_4)$ as well (since
$(P_5)$ for $5$ odd gives $v_5^*(K_5) = 0$, and the dim and
$K_4^{\mathrm{embed}} \subset K_5$ force equality).  So $\xi^{(5)} =
v_3 - \gamma_2 v_4$ (taken modulo $\mathbb{Q}(q) v_2 = K_3^{\mathrm{embed}}$),
giving $\mu_5 = -\gamma_2/1 = -\gamma_2$.
\end{proof}

\begin{lemma}[Recursion for $\mu_l$, $l$ odd $\ge 7$]\label{lem:mu-rec}
Assume $(R_{l-1})$ holds at level $l-1$ even.  Then
\begin{equation}\label{eq:mu-rec}
   \mu_l \;=\; \frac{-\gamma_{l-3}}{1 - \gamma_{l-4}\, \nu_{l-3}}.
\end{equation}
\end{lemma}

\begin{proof}
Pick $\xi^{(l)} \in K_l$ new generator. By the May-9 reduction
analysis at level $l$: $R_l'(\xi^{(l)}) = y$ lies in
$D_l \cap K_{l-1}^{\mathrm{embed}} = K_{l-1}^{\mathrm{embed}} \cap
\mathrm{span}(v_2,\ldots,v_{l-2})$.  By $(P_{l-1})$ for $l-1$ even,
$v_{l-1}^*(K_{l-1}) \neq 0$ but we additionally need to know the
structure of $K_{l-1}\cap\{v_{l-1}^*=0\}$.

By $(R_{l-1})$ assumed: $v_{l-2}^*$ and $v_{l-1}^*$ are proportional
on $K_{l-1}$ with nonzero ratio; hence
$K_{l-1} \cap \{v_{l-1}^*=0\} = K_{l-1}\cap\{v_{l-2}^*=0\}$, and by
nesting $+$ dimension this equals $K_{l-3}^{\mathrm{embed}}$
(in $V_{(l-2,1)}$).

So $y \in K_{l-3}^{\mathrm{embed}}$, which has support on
$v_2,\ldots,v_{l-3}$ since $K_{l-3} \subset V_{(l-4,1)}$
(natural embedding).  In particular, $y_{v_{l-3}}$ may be nonzero
while $y_{v_{l-2}} = 0$.

By Strong Image at level $n = l$:
\begin{align*}
\xi^{(l)}_{v_{l-1}} &= \frac{y_{v_{l-2}} - \gamma_{l-3} y_{v_{l-3}}}{(1+q)^{l-2}}
                   \;=\; \frac{-\gamma_{l-3}\, y_{v_{l-3}}}{(1+q)^{l-2}},\\
\xi^{(l)}_{v_{l-2}} &= \frac{y_{v_{l-3}} - \gamma_{l-4}\, y_{v_{l-4}}}{(1+q)^{l-2}}.
\end{align*}
Hence
\[
   \mu_l \;=\; \frac{\xi^{(l)}_{v_{l-1}}}{\xi^{(l)}_{v_{l-2}}}
   \;=\; \frac{-\gamma_{l-3}\, y_{v_{l-3}}}{y_{v_{l-3}} - \gamma_{l-4}\, y_{v_{l-4}}}
   \;=\; \frac{-\gamma_{l-3}}{1 - \gamma_{l-4}\,(y_{v_{l-4}}/y_{v_{l-3}})}.
\]

Finally the ratio $y_{v_{l-4}}/y_{v_{l-3}}$ on $K_{l-3}^{\mathrm{embed}}$:
since $K_{l-3} = K_{l-5}^{\mathrm{embed}} \oplus \mathbb{Q}(q)\eta^{(l-3)}$
and $K_{l-5}^{\mathrm{embed}} \subset \mathrm{span}(v_2,\ldots,v_{l-5})$
(both $v_{l-4}, v_{l-3}$ coordinates vanish on $K_{l-5}^{\mathrm{embed}}$),
the ratio equals $\eta^{(l-3)}_{v_{l-4}}/\eta^{(l-3)}_{v_{l-3}} =
\nu_{l-3}$.  Substituting gives \eqref{eq:mu-rec}.
\end{proof}

\begin{lemma}[Recursion for $\nu_l$, $l$ even $\ge 6$]\label{lem:nu-rec}
\begin{equation}\label{eq:nu-rec}
   \nu_l \;=\; \frac{1 - \gamma_{l-3}/\mu_{l-1}}{(1+q)^{l-2}\, C_l},
   \qquad
   C_l \;=\; \frac{\beta_{l-2}\beta'_{l-1}}{(1+q)^{l-1}}.
\end{equation}
\end{lemma}

\begin{proof}
Pick $\eta^{(l)} \in K_l$ new generator.  Then $R_l'(\eta^{(l)}) =
z \in K_{l-1}^{\mathrm{embed}}$ where $K_{l-1}$ is at level $l-1$
odd, so $z$ has support on $\mathrm{span}(v_2,\ldots,v_{l-2})$ (by
$(P_{l-1})$).  By Image Lemma: $\eta^{(l)}_{v_l} = C_l z_{v_{l-2}}$.
Since $\eta^{(l)} \notin K_{l-2}^{\mathrm{embed}}$ implies
$\eta^{(l)}_{v_l} \neq 0$, we have $z_{v_{l-2}} \neq 0$.

By Strong Image:
$\eta^{(l)}_{v_{l-1}} = (z_{v_{l-2}} - \gamma_{l-3} z_{v_{l-3}})/(1+q)^{l-2}$.
So
\[
\nu_l = \frac{\eta^{(l)}_{v_{l-1}}}{\eta^{(l)}_{v_l}}
      = \frac{z_{v_{l-2}} - \gamma_{l-3} z_{v_{l-3}}}{(1+q)^{l-2}\, C_l\, z_{v_{l-2}}}
      = \frac{1 - \gamma_{l-3}\,(z_{v_{l-3}}/z_{v_{l-2}})}{(1+q)^{l-2}\, C_l}.
\]

The ratio $z_{v_{l-3}}/z_{v_{l-2}}$ on $K_{l-1}$ (odd): since
$K_{l-1} = K_{l-3}^{\mathrm{embed}} \oplus \mathbb{Q}(q)\xi^{(l-1)}$
and $K_{l-3}^{\mathrm{embed}} \subset \mathrm{span}(v_2,\ldots,v_{l-4})$
(the latter inclusion since $K_{l-3} \subset V_{(l-4,1)}$ embeds as
$\mathrm{span}(v_2,\ldots,v_{l-4})$ in $V_{(l-2,1)}$, with no
$v_{l-3}, v_{l-2}$ components), the ratio is determined by $\xi^{(l-1)}$:
$z_{v_{l-3}}/z_{v_{l-2}} = \xi^{(l-1)}_{v_{l-3}}/\xi^{(l-1)}_{v_{l-2}} = 1/\mu_{l-1}$.

Substituting yields \eqref{eq:nu-rec}.
\end{proof}

\section{Sign induction}\label{sec:signs}

We now prove an alternating sign property of the slopes at any
positive real $q$, treating $q > 0$ as a real variable in $\mathbb{R}_{>0}$.

\begin{lemma}[Positivity of Hoefsmit constants]\label{lem:positivity}
At any real $q > 0$, the constants $\beta_i, p_i, q_i, \gamma_i$
are strictly positive, and $\beta'_i$ is strictly negative.  In
particular, $C_l = \beta_{l-2} \beta'_{l-1}/(1+q)^{l-1} < 0$.
\end{lemma}

\begin{proof}
$[k] = (q^k-1)/(q-1) > 0$ at $q > 0, q \neq 1$, and $[k]|_{q=1} = k > 0$.
Therefore $\beta_i = q[i-1]/[i] > 0$, $p_i = [i+1]/([2][i]) > 0$,
$q_i = 1/[2] > 0$, and $\gamma_i = \beta_i p_{i+1} > 0$.
$\beta'_i = -[i+1]/[i] < 0$.  Hence $C_l = (\text{positive})\cdot(\text{negative})/(\text{positive}) < 0$.
\end{proof}

\begin{theorem}[Alternating sign]\label{thm:alt-sign}
For all real $q > 0$:
\begin{itemize}[topsep=0.3em,itemsep=0.2em,leftmargin=2em]
\item $\nu_l < 0$ for every even $l \ge 4$.
\item $\mu_l < 0$ for every odd $l \ge 5$.
\end{itemize}
The proof is by strong induction simultaneously establishing
$(P_l), (Q_l), (R_l)$ along with these signs.
\end{theorem}

\begin{proof}
Order the levels $4 < 5 < 6 < 7 < \cdots$ and induct on $l$.

\textbf{Bases.}
\begin{itemize}[topsep=0.3em,itemsep=0.2em,leftmargin=2em]
\item Level $3$: $K_3 = \mathbb{Q}(q) v_2$, $(P_3), (Q_3)$ trivial.
\item Level $4$: $K_4 = \mathrm{span}(v_2, v_3 - \gamma_2 v_4)$
  (computed). $(P_4), (Q_4), (R_4)$ all hold.  $\nu_4 = -1/\gamma_2 < 0$
  by Lemma~\ref{lem:positivity}.
\item Level $5$: $K_5 = \mathrm{span}(v_2, v_3 - \gamma_2 v_4)$
  (extended by zero on $v_5$).  $(P_5), (Q_5)$ hold, $\mu_5 = -\gamma_2 < 0$.
\end{itemize}

\textbf{Inductive step.} Suppose all properties (sign, $(P), (Q),
(R)$) hold at all levels $< l$.

\paragraph{Case $l$ odd, $l \ge 7$.}
By induction, $(R_{l-1})$ holds (since $l-1$ even is $< l$, and
$(R_{l-1})$ follows from $(P_{l-1}), (Q_{l-1})$ via May-9
Reduction Theorem).  Apply Lemma~\ref{lem:mu-rec}:
$\mu_l = -\gamma_{l-3}/(1 - \gamma_{l-4}\nu_{l-3})$.

By induction, $\nu_{l-3} < 0$ (since $l-3$ even is $< l$).  By
Lemma~\ref{lem:positivity}, $\gamma_{l-4} > 0$.  Hence $\gamma_{l-4}
\nu_{l-3} < 0$, so $1 - \gamma_{l-4}\nu_{l-3} > 1 > 0$.  Therefore
\[
   \mu_l \;=\; \frac{-\gamma_{l-3}}{\text{positive}} \;<\; 0.
\]

We need to check that $\xi^{(l)}_{v_{l-2}} \neq 0$ (so that $\mu_l$
is finite).  By Lemma~\ref{lem:mu-rec}'s computation,
$\xi^{(l)}_{v_{l-2}} = (y_{v_{l-3}}/(1+q)^{l-2})(1 - \gamma_{l-4}\nu_{l-3})$,
which is nonzero since both factors are.

Properties $(P_l), (Q_l)$ (with $(Q_l)$ for $l$ odd automatic):
established via the May-8 evening induction from $(P_{l'}), (R_{l'})$
at smaller $l'$.

\paragraph{Case $l$ even, $l \ge 6$.}
By induction, $\mu_{l-1} < 0$ (since $l-1$ odd is $< l$, with
$(P_{l-1}), (Q_{l-1})$ valid).  Apply Lemma~\ref{lem:nu-rec}:
$\nu_l = (1 - \gamma_{l-3}/\mu_{l-1})/((1+q)^{l-2} C_l)$.

By induction, $\mu_{l-1} < 0$, $\gamma_{l-3} > 0$, hence $\gamma_{l-3}
/\mu_{l-1} < 0$, so $1 - \gamma_{l-3}/\mu_{l-1} > 1 > 0$.  By
Lemma~\ref{lem:positivity}, $C_l < 0$, so $(1+q)^{l-2} C_l < 0$.
Therefore
\[
   \nu_l \;=\; \frac{\text{positive}}{\text{negative}} \;<\; 0.
\]

For $(P_l)$ at level $l$ even: by May-8 evening induction,
$(P_l) \Leftarrow (P_{l-1}), (Q_{l-1})$ (which we have).

For $(Q_l)$ at level $l$ even: see Theorem~\ref{thm:Q-main} below.

For $(R_l)$ at level $l$ even: by May-9 Reduction Theorem,
$(R_l) \Leftarrow (P_l), (Q_l), \text{Nesting}$, all of which we
have.
\end{proof}

\section{The main theorem}\label{sec:main}

\begin{theorem}[Property $(Q_m)$ for all even $m$]\label{thm:Q-main}
For every even $m \ge 4$, the linear functional $v_{m-1}^*$ is not
identically zero on $K_m$.  Equivalently, there exists $w \in K_m$
with $w_{v_{m-1}} \neq 0$.
\end{theorem}

\begin{proof}
The case $m = 4$ is base: $K_4 = \mathrm{span}(v_2, v_3 - \gamma_2 v_4)$
contains $v_3 - \gamma_2 v_4$ with $v_3$-coordinate $1 \neq 0$.

For $m \ge 6$ even: by Theorem~\ref{thm:alt-sign} (sign induction
established at all levels $< m$), $\mu_{m-1} < 0$ at every $q > 0$.
By Lemma~\ref{lem:positivity}, $\gamma_{m-3} > 0$.  Therefore
\[
   \mu_{m-1} \;-\; \gamma_{m-3} \;<\; 0
\]
at every $q > 0$.  Since $\mu_{m-1} - \gamma_{m-3}$ is a rational
function in $q$ that is strictly negative at every positive real $q$,
it cannot be identically zero.  Hence $\mu_{m-1} - \gamma_{m-3} \neq 0$
in $\mathbb{Q}(q)$.

Pick $\xi^{(m-1)} \in K_{m-1}$ a new generator of $K_{m-1}$ with
$\xi^{(m-1)}_{v_{m-2}} \neq 0$ (exists by $(Q_{m-1})$, automatic for
$m-1$ odd).  By Theorem~\ref{thm:alt-sign}'s argument,
$\xi^{(m-1)}_{v_{m-3}} \neq 0$ as well.  Embed $\xi^{(m-1)}$ into
$V_{(m-1,1)}$ via the natural embedding, and consider the linear
functional
\[
   \phi_m \;:=\; v_{m-2}^* - \gamma_{m-3}\, v_{m-3}^*
\]
applied to $\xi^{(m-1)}$:
\[
   \phi_m(\xi^{(m-1)}) \;=\; \xi^{(m-1)}_{v_{m-2}} - \gamma_{m-3}\, \xi^{(m-1)}_{v_{m-3}}
   \;=\; \xi^{(m-1)}_{v_{m-3}}\bigl(\mu_{m-1} - \gamma_{m-3}\bigr) \;\neq\; 0.
\]

Now lift $\xi^{(m-1)}$ to $w \in K_m$ via $R_m'^{-1}$: pick any
$w \in V_{(m-1,1)}$ with $R_m'(w) = \xi^{(m-1)}$.  By the May-8
evening Reduction Theorem (using the Image Lemma and $K_{m-1}$
analysis), such $w$ lies in $K_m$ provided $\xi^{(m-1)} \in
D_m \cap K_{m-1}^{\mathrm{embed}}$.  Since $\xi^{(m-1)}_{v_{m-1}} = 0$
($(P_{m-1})$ for $m-1$ odd), $\xi^{(m-1)} \in
\mathrm{span}(v_2,\ldots,v_{m-2}) \subset D_m$.  And by the
inductive hypothesis $\xi^{(m-1)} \in K_{m-1}^{\mathrm{embed}}$.  So
$w \in K_m$.

By Strong Image at level $n = m$ applied to $w$ with $y = \xi^{(m-1)}$:
\[
   w_{v_{m-1}} \;=\; \frac{\xi^{(m-1)}_{v_{m-2}} - \gamma_{m-3}\,\xi^{(m-1)}_{v_{m-3}}}{(1+q)^{m-2}}
   \;=\; \frac{\phi_m(\xi^{(m-1)})}{(1+q)^{m-2}} \;\neq\; 0.
\]
This proves $(Q_m)$.
\end{proof}

\begin{corollary}[Hook rank formula, unconditional]\label{cor:hook-rank}
For every $n \ge 2$,
\[
   \mathrm{rank}\,\Pi^{S_n}|_{V_{(n-1,1)}} \;=\; \big\lceil \tfrac{n-2}{2} \big\rceil.
\]
\end{corollary}

\begin{proof}
By Theorem~\ref{thm:Q-main}, $(Q_m)$ holds for every even $m \ge 4$.
Combined with the May-9 Reduction Theorem and the rank-formula
proof in May-9 \S4, the rank formula holds for all $n \ge 2$.
\end{proof}

\section{Computational verification}\label{sec:numerical}

The Strong Image Coefficient Formula \eqref{eq:strong-image} is
verified symbolically at $n = 4, 5, 6, 7, 8$ in
\texttt{verify\_strong\_image.py}, where the formula is
matched against direct linear-algebra computation of $w$ given
$y$.

The slopes $\mu_l, \nu_l$ have been verified explicitly:
\begin{itemize}[topsep=0.3em,itemsep=0.2em,leftmargin=2em]
\item $\nu_4 = -1/\gamma_2 = -(q+1)(q^2+q+1)/(q(q^2+1))$.
\item $\mu_5 = -\gamma_2 = -q(q^2+1)/((q+1)(q^2+q+1))$.
\item $\nu_6 = (\text{computed in May-8 evening Table})$.
\item $\mu_7 = -\gamma_2\gamma_4/(\gamma_2+\gamma_3) = -q(q^2+1)(q^2-q+1)(q^2+q+1)/((q+1)(2q^2-q+2)(q^4+q^3+q^2+q+1))$.
\end{itemize}

The identity $\mu_{m-1} - \gamma_{m-3}$ for $(Q_m)$ at $m = 6, 8$:
\begin{itemize}[topsep=0.3em,itemsep=0.2em,leftmargin=2em]
\item $m=6$: $\mu_5 - \gamma_3 = -(\gamma_2 + \gamma_3)
  = -q(2q^2-q+2)/((q+1)(q^2+1))$, nonzero.
\item $m=8$: $\mu_7 - \gamma_5 = -q(q^2+1)(3q^4-2q^3+4q^2-2q+3)/((q+1)(q^2-q+1)(q^2+q+1)(2q^2-q+2))$,
  nonzero.
\end{itemize}

Both negative at $q > 0$ as predicted by Theorem~\ref{thm:alt-sign}.

\section{Honest assessment}\label{sec:honest}

\paragraph{What this paper proves.}
\begin{enumerate}[topsep=0.3em,itemsep=0.2em,leftmargin=2em]
\item Theorem~\ref{thm:strong-image} (Strong Image Coefficient
  Formula).  Rigorous via direct triangular matrix inversion of
  $R_{n-1}''$ on $V_{(n-1,1)}$.
\item Lemmas~\ref{lem:mu-rec} and \ref{lem:nu-rec}: recursive
  formulas for the slopes $\mu_l, \nu_l$.
\item Theorem~\ref{thm:alt-sign} (alternating sign): both slopes
  strictly negative at every positive real $q$.
\item Theorem~\ref{thm:Q-main}: $(Q_m)$ holds for all even $m \ge 4$.
\item Corollary~\ref{cor:hook-rank}: hook rank formula proved
  unconditionally.
\end{enumerate}

\paragraph{Where Hoefsmit-specific computation enters.}
The proof crucially uses positivity of the explicit Hoefsmit
constants $\beta_i, p_i, \gamma_i$ at $q > 0$, and negativity of
$\beta'_i$.  These hold because $V_{(n-1,1)}$'s seminormal block
matrices have a uniform sign structure under our normalization.

\paragraph{Inductive scope.}
The induction proceeds simultaneously on $(P_l), (Q_l), (R_l)$ for
$l \le m$, with the sign properties of $\mu_l, \nu_l$ as additional
inductive invariants.  Each level requires properties at all
strictly smaller levels.  No level uses its own $(R_l)$ or $(Q_l)$
in deriving the next level.

\paragraph{Files.}
\begin{itemize}[topsep=0.3em,itemsep=0.2em,leftmargin=2em]
\item \texttt{\textasciitilde/projects/proofs/2026-05-10-Q-m-via-strong-image.tex} (this writeup).
\item \texttt{\textasciitilde/projects/scratch/2026-05-09-nesting-identity/check\_identity.py}
  (Hoefsmit constant values).
\item \texttt{\textasciitilde/projects/scratch/2026-05-08-kernel-dichotomy/check\_n8\_proportionality.py}
  (computational verification of $(Q_m), (R_m)$ at $m \le 8$).
\end{itemize}

\paragraph{What remains.}
With $(Q_m)$ proved, the May-9 Reduction Theorem closes the hook rank
formula unconditionally.  Future directions:
\begin{itemize}[topsep=0.3em,itemsep=0.2em,leftmargin=2em]
\item Closed form for the $\mu_l, \nu_l$ progressions (continued
  fraction in the Hoefsmit gammas).
\item Generalization to other 2-row partitions $(n-r, r)$ for $r \ge 2$.
\item Cellular interpretation: does the alternating sign reflect a
  duality between the seminormal form and another preferred basis
  (Kazhdan-Lusztig, web bases, etc.)?
\end{itemize}

\end{document}
